\section{MuJoCo 기반 교육용 실험실 설계}

\subsection{교육 목표}

이 실습 환경이 목표로 하는 것은 고성능 로봇 학습 벤치마크가 아니라, 제어 개념을 순서대로 확장할 수 있는 로컬 실험 환경이다.
로봇 시뮬레이션은 연구와 교육 모두에서 유용하지만, 시뮬레이터의 목적과 한계를 명확히 해야 한다 \cite{choi2021simulationrobotics,collins2021physicsreview}.

제어 교육에서 실험실의 역할은 단순히 장비를 보여주는 것이 아니라, 학생이 같은 도구 체계 안에서 모델링, 제어기 조정, 실행 결과 해석을 반복하게 만드는 데 있다.
기존 제어 교육 연구도 표준화된 소프트웨어 환경과 공유 가능한 실험 자원이 학부 실험실의 운영 부담을 줄이고 학습 내용을 반복 가능하게 만들 수 있음을 강조했다 \cite{dixon2002matlablab}.
로봇 조작 교육에서도 원격 또는 가상 실험실을 통해 학생이 로봇 매니퓰레이터 제어 알고리즘을 안전하게 시험하고, 결과 그래프와 움직임을 함께 보며 파라미터 감각을 익히게 하려는 시도가 있었다 \cite{sartorius2006virtualremote,hoenig2016seamless}.
여기서 핵심은 여러 과제의 성능 순위를 매기는 비교보다, 학생이 어떤 순서로 개념을 익히는가이다.
각 랩은 하나의 개념을 드러내는 작은 실험으로 만들고, 학생이 YAML 설정 파일을 바꾸고 로그와 플롯을 보며 원인과 결과를 직접 연결하도록 설계한다.

실험실 설계 원칙은 다음과 같다.
\begin{itemize}
  \item 모든 랩은 MuJoCo 물리 엔진을 사용한다.
  \item 모든 실험은 YAML 설정 파일로 재현 가능해야 한다.
  \item 실행 결과는 CSV, NPZ, summary JSON, 플롯 이미지로 저장한다.
  \item 뷰어 기반 실시간 데모와 창 없는 자동 실행을 모두 지원한다.
  \item 수식, 코드, 그래프가 서로 대응되도록 작고 읽기 쉬운 구현을 유지한다.
\end{itemize}

여기서 파일 형식도 학습 흐름의 일부이다.
MuJoCo는 관절, 물체, 접촉, 중력 같은 물리 현상을 계산해주는 물리 시뮬레이터이다.
여기서 산출물은 한 번의 실행 뒤 자동으로 남는 파일 묶음을 뜻한다.
YAML 설정 파일은 사람이 읽기 쉬운 실험 레시피이다.
예를 들어 질량, 강성, 감쇠, 목표 위치, 실행 시간을 숫자로 기록한다.
CSV 로그는 시간에 따라 변한 값들을 표처럼 저장한 파일이다.
위치, 속도, 힘, 에너지 같은 열을 엑셀이나 파이썬으로 다시 열어볼 수 있다.
NPZ는 NumPy 배열을 보존하기 위한 압축 파일이고, summary JSON은 한 번의 실험을 대표하는 핵심 숫자를 담는다.
플롯 이미지는 강의 자료나 보고서에서 바로 볼 수 있는 그래프이다.
MuJoCo 뷰어는 실시간으로 로봇과 물리 세계를 눈으로 확인하는 창이고, 창 없는 실행은 화면을 띄우지 않고 자동으로 실험을 돌리는 방식이다.
뷰어에서는 움직임을 보며 직관을 얻고, 창 없는 실행은 같은 실험을 다시 확인할 로그와 플롯을 남긴다.
effort 지표는 실제 전력 소모를 직접 잰 값이 아니라, 제어기나 actuator가 얼마나 세게 밀거나 돌리려 했는지를 비교하기 위한 대리 지표이다.

\begin{table}[H]
  \centering
  \caption{저장된 산출물에서 이론으로 돌아가는 길}
  \label{tab:artifact-theory-bridge}
  \small
  \setlength{\tabcolsep}{4pt}
  \renewcommand{\arraystretch}{1.12}
  \begin{tabularx}{\linewidth}{@{}>{\raggedright\arraybackslash}p{0.21\linewidth}>{\raggedright\arraybackslash}p{0.33\linewidth}>{\raggedright\arraybackslash}X@{}}
    \toprule
    산출물 & 열어볼 것 & 다시 물어볼 질문 \\
    \midrule
    \texttt{config.yaml} & 이번 실행에서 바꾼 질량, 강성, 감쇠, PID 이득, 목표 위치, 벽 위치, 실행 시간 & 이 값은 앞 절 수식의 어느 항을 건드렸는가? \\
    \texttt{log.csv} & 시간에 따른 위치, 속도, 오차, 제어 입력, effort 지표, 침투 깊이 같은 로그 열 & 입력이나 목표가 바뀐 뒤 응답은 어떤 순서로 따라왔는가? \\
    \texttt{states.npz} & 숫자 열을 NumPy 배열로 다시 불러와 계산할 원자료 & 기본 플롯 말고도 에너지, 피크, 평균 effort처럼 다시 계산해 볼 신호가 있는가? \\
    \texttt{summary.json} & 오버슈트, 정착시간, 최대 effort, 최종 오차, 최대 침투 깊이 같은 대표 숫자 & 눈으로 본 차이가 요약 숫자로도 같은 방향을 가리키는가? \\
    \texttt{plots/*.png} & \texttt{--plot}을 사용했을 때 저장되는 위치, 속도, 힘 또는 effort의 시간 그래프 & 흔들림, 느린 수렴, 포화, 접촉 시작은 그래프의 어느 구간에서 보이는가? \\
    \texttt{report.html}, \texttt{worksheet.md} & 실행 전 예측, 실행 후 관찰, 핵심 플롯, 다음 실험 메모 & 예상과 달랐다면 다음에는 무엇 하나만 바꿔 볼까? \\
    인터랙티브 산출물 & \texttt{interaction\_events.json}, \texttt{learner\_snapshot.json}, \texttt{learner\_tuned\_config.yaml}처럼 조작 기록이나 재실행 설정이 남은 파일 & 실행 중 손으로 바꾼 값이 다음 재현 가능한 가설로 바뀌었는가? \\
    \bottomrule
  \end{tabularx}
\end{table}

표~\ref{tab:artifact-theory-bridge}의 파일들은 실행 뒤에 남는 부속물이 아니라, 실험을 다시 읽기 위한 단서이다.

\texttt{config.yaml}에는 실행 전 가설이, \texttt{log.csv}와 저장된 플롯에는 그 가설이 시간응답으로 나타난 모습이 남는다.
\texttt{summary.json}은 긴 그래프를 몇 개의 숫자로 접어 보는 도구이고, \texttt{worksheet.md}는 다음 실험을 고르는 메모장에 가깝다.
일부 산출물은 실행 방식에 따라 달라진다.
플롯은 \texttt{--plot}을 사용했을 때 저장되고, 인터랙티브 산출물은 학습자가 슬라이더나 버튼으로 값을 바꾸고 관찰을 남겼을 때 의미가 커진다.

다만 이 산출물은 교육용 시뮬레이션을 해석하기 위한 자료이다.
Lab04의 \path|tau_cmd_*|와 \path|current_proxy_*|는 MuJoCo \path|data.actuator_force|에서 온 시뮬레이션 effort 지표로 읽는다.
\path|force_virtual_0|, \path|force_virtual_spring_0|, \path|force_virtual_damping_0|은 설정한 가상 벽 모델에서 계산한 신호로 읽는다.
즉, 이 값들은 실제 모터 전류나 하드웨어 힘 센서값이라기보다, 같은 설정 안에서 어느 쪽이 더 큰 요구와 접촉 반응을 만들었는지 비교하는 눈금이다.

\subsection{랩 결과를 읽는 순서}

랩 결과는 로봇이 얼마나 멋지게 움직였는지보다, 설정값 변화가 시간 그래프에 어떻게 찍혔는지부터 읽는다.
앞선 이론 장들에서 배운 수식이 시간 그래프에서 어떤 변화로 나타나는지 확인하기 위한 관찰 장치이기 때문이다.
결과를 볼 때는 다음 순서를 따라가는 편이 좋다.

\begin{enumerate}
  \item 1단계에서는 설정 파일에서 무엇을 바꾸었는지 확인한다. 질량, 강성, 감쇠, 목표 위치, 벽 위치, 실행 시간이 여기에 해당한다.
  \item 그다음에는 입력과 목표를 본다. 외력이 들어갔는지, 목표 위치가 계단처럼 바뀌었는지, 부드러운 궤적인지 확인한다.
  \item 이제 위치와 속도를 본다. 최종 위치 오차, 오버슈트, 진동, 정착시간이 여기서 보인다.
  \item 마지막에는 힘, 제어 입력, 에너지를 본다. 같은 위치 응답이라도 더 큰 힘을 썼는지, 에너지가 저장되었다가 사라지는지, 포화에 가까운지 확인한다.
\end{enumerate}

이 순서가 중요한 이유는 임피던스가 하나의 숫자가 아니라 힘과 운동의 관계이기 때문이다.
즉, 위치 그래프만 따로 보는 것이 아니라 같은 시간의 힘, 제어 입력, effort 그래프를 나란히 보며 얼마나 밀었을 때 얼마나 움직였는지를 읽는다.
강성을 올리면 정상상태 위치 오차가 줄어드는 경향이 있지만, 동시에 같은 오차에서 더 큰 힘과 더 큰 가상 퍼텐셜 에너지를 만든다.
감쇠를 올리면 흔들림과 오버슈트가 줄어드는 경향이 있지만, 너무 크면 움직임이 둔해지고 목표에 늦게 도달할 수 있다.
질량 또는 유효 질량이 커지면 같은 힘에 대한 가속도가 작아져 천천히 반응하고, 고유진동수도 달라진다.
그래서 한 그래프만 보면 안 된다.
위치 오차가 작아졌다는 사실과 제어 입력이 커졌다는 사실을 함께 봐야, 그 파라미터가 실제로 어떤 대가를 치렀는지 이해할 수 있다.
그래프를 볼 때는 먼저 세 가지를 묻는다.
최종값은 어디에 멈췄는가?
멈추기 전 얼마나 흔들렸는가?
그 과정에서 힘이나 제어 입력은 얼마나 커졌는가?
위치/속도 플롯을 읽기 전에 먼저 이 실행이 정상상태까지 갔는지 확인한다.
그래프의 마지막 점 하나를 최종값으로 삼으면, 아직 내려가거나 흔들리는 과도응답을 정상상태 오차로 착각할 수 있다.
꼬리 구간에서 위치 곡선이 거의 평평하고 속도가 0 근처로 줄었을 때만 그 값을 최종 위치로 잡는다.
아직 기울기가 남아 있거나 피크가 계속 줄어드는 중이라면, 먼저 실행 시간을 늘리거나 \texttt{summary.json}의 정착시간과 함께 확인한다.
오버슈트는 이렇게 잡은 최종값을 기준으로 가장 큰 피크를 재고, 정착시간은 그 최종값 주변의 허용 띠 안에 들어간 뒤 다시 벗어나지 않는 첫 시각으로 읽는다.
허용 띠는 최종값 주변에서 이 정도 오차까지는 충분히 도착했다고 보겠다고 정한 범위이다.
예를 들어 최종값의 $\pm2\%$ 안쪽을 허용 띠로 둘 수 있다.
강성은 주로 최종 위치 오차와 힘-변위 기울기를 바꾸고, 감쇠는 오버슈트와 정착시간을 바꾸며, 유효 질량은 가속도와 진동 주기를 바꾼다.
따라서 각 랩의 플롯은 보기 좋은 움직임만 확인하는 그림이 아니라, 앞에서 정리한
\begin{equation}
  |e_{\mathrm{ss}}|\approx \frac{|F|}{k},
  \qquad
  \omega_n=\sqrt{\frac{k}{m}},
  \qquad
  \zeta=\frac{d}{2\sqrt{mk}}
\end{equation}
같은 관계가 시간응답에 나타난 변화를 찾는 자료이다.
여기서 $|e_{\mathrm{ss}}|$는 위치 플롯 끝에 남은 오차, $\omega_n$은 진동 주기의 짧고 김, $\zeta$는 피크가 얼마나 빨리 줄어드는지로 읽는다.
다만 이 식들은 모든 랩에 같은 강도로 적용되지 않는다.
Lab01/02처럼 1축 등가 질량을 명확히 둘 수 있는 랩에서는 비교적 직접 대응시켜 읽고, Lab03/04에서는 정확한 스칼라 공식보다 변화 방향을 읽는 출발점으로 사용한다.
또 논문 기호와 설정 파일 이름이 항상 같지는 않다.
Lab01은 \texttt{stiffness}와 \texttt{damping}, Lab02는 PID의 \texttt{kp}, \texttt{ki}, \texttt{kd}, Lab03은 \texttt{task\_kp}, \texttt{task\_kd}와 DLS 설정, Lab04 가상 벽은 \texttt{virtual\_wall.stiffness}와 \texttt{virtual\_wall.damping}을 만진다.
특히 Lab04에서 DLS damping은 역기구학을 안정화하는 수치 완화이고, \texttt{virtual\_wall.damping}은 접근 속도에 작용하는 교육용 벽 모델 계수이지 검증된 끝단 물리 감쇠비가 아니다.

\begin{table}[!htbp]
  \centering
  \caption{랩에서 먼저 봐야 할 신호와 연결되는 이론 질문}
  \label{tab:lab-observation-guide}
  \small
  \renewcommand{\arraystretch}{1.12}
  \begin{tabularx}{\linewidth}{@{}p{0.12\linewidth}p{0.22\linewidth}p{0.25\linewidth}X@{}}
    \toprule
    랩 & 바꾸는 값 & 먼저 볼 신호 & 붙잡을 질문 \\
    \midrule
    Lab01 & $m$, $b$, $k$, 초기 위치, 외력 & 위치, 속도, 가속도 추정, 운동에너지, 퍼텐셜 에너지 & 고유진동수와 감쇠비가 오버슈트와 정착시간을 어떻게 바꾸는가? \\
    Lab02 & P, I, D 이득, 포화, 지연, 노이즈 & 목표-실제 위치, 오차, 제어 입력, 오버슈트, 정상상태 오차 & P 항은 왜 가상 스프링처럼 보이고, D 항은 왜 가상 댐퍼처럼 보이는가? \\
    Lab03 & 관절 목표, 작업공간 목표, 자코비안 조건, 감쇠 최소제곱 계수 & 관절각, 끝단 위치, 토크 명령 또는 전류 프록시 형태의 actuator effort 지표, 조작성, 조건수 & 같은 손끝 목표라도 자세와 특이점 근처에서 왜 필요한 관절 움직임과 effort가 달라지는가? \\
    Lab04 & 작업공간 목표, 강성, 감쇠, 벽 위치, 목표 후퇴 설정 & 끝단 오차, 속도, 침투 깊이, 계산된 가상 벽 힘, 액추에이터 힘 지표 & 강성과 감쇠를 바꾸면 접촉 순간의 힘, 튕김, 침투 깊이, 목표 후퇴가 어떻게 달라지는가? \\
    \bottomrule
  \end{tabularx}
\end{table}

표~\ref{tab:lab-observation-guide}의 질문은 정답을 맞히기 위한 목록이 아니라, 플롯을 읽기 위한 순서이다.
예를 들어 Lab01에서 강성을 크게 하면 고유진동수가 커지고, 같은 변위에 저장되는 퍼텐셜 에너지도 커진다.
그 결과 빠르게 움직이지만 감쇠가 충분하지 않으면 더 잘 흔들릴 수 있다.
Lab04의 가상 벽에서도 같은 생각이 반복된다.
벽 강성이 크면 같은 침투 깊이에서 계산된 벽 힘의 기울기는 커진다.
침투 깊이는 같은 목표, 후퇴 게인, 포화 조건, DLS 조건에서 작아지는 경향을 보일 수 있지만 항상 단조롭게 줄어든다고 단정하면 안 된다.
따라서 안전한 접촉을 보려면 침투 깊이만이 아니라 벽 힘, 속도, 목표 후퇴량, 제어 입력을 함께 봐야 한다.

개별 랩으로 들어가기 전 공통 루틴을 정하자.
실행 전에는 이번에 바꾼 파라미터가 어느 플롯을 어떻게 바꿀 것인지 한 문장으로 예측한다.
실행 후에는 설정값, 위치와 속도, 힘 또는 effort, 에너지 또는 침투 깊이를 같은 순서로 확인한다.
마지막에는 예상과 맞은 점, 달랐던 점, 다음에 바꿔 볼 파라미터를 짧게 기록한다.
각 랩의 산출물은 실행 성공 여부만이 아니라 해석 기록이다.
예를 들어 실행 전에는 Lab01의 \texttt{stiffness}를 키우면 최종 위치 오차는 줄지만 effort 피크는 커질 것이다라고 적는다.
Lab04 가상 벽에서는 \texttt{virtual\_wall.stiffness}를 키우면 같은 침투에서 계산된 벽 힘은 커지고, 같은 나머지 조건에서는 침투 깊이가 줄어드는 경향이 있는지 확인하겠다고 적는 편이 더 정확하다.
실행 후에는 실제 플롯에서 위치 오차가 줄었는지, 속도나 오버슈트가 커졌는지, effort가 한계에 가까워졌는지를 같은 문장으로 되짚는다.
이렇게 예측 한 줄, 실행, 플롯 해석 한 줄을 반복하면 수식이 암기 항목이 아니라 실험을 읽는 도구가 된다.

\begin{table}[!htbp]
  \centering
  \caption{랩별 플롯 확인 메모}
  \label{tab:lab-first-plot-check}
  \small
  \setlength{\tabcolsep}{4pt}
  \renewcommand{\arraystretch}{1.12}
  \begin{tabularx}{\linewidth}{@{}>{\raggedright\arraybackslash}p{0.12\linewidth}>{\raggedright\arraybackslash}p{0.28\linewidth}>{\raggedright\arraybackslash}p{0.28\linewidth}>{\raggedright\arraybackslash}X@{}}
    \toprule
    랩 & 기대 모습 & 첫 이상 신호 & 다음 확인 \\
    \midrule
    Lab01 & 위치가 평형점으로 돌아가고, 부족감쇠에서는 줄어드는 진동, 과감쇠에서는 느린 수렴이 보인다 & 진동이 줄지 않거나, 힘 입력이 끝난 뒤에도 에너지가 계속 커진다 & 외력 종류, 실행 시간, $m$, $b$, $k$, 감쇠 부호, 에너지 감소 여부를 먼저 본다 \\
    Lab02 & 목표와 실제 위치의 차이가 줄고, 제어 입력이 포화 한계에 오래 붙어 있지 않는다 & 입력이 한계에서 평평하게 잘리거나, 오차가 줄지 않고 커진다 & 피드백 부호, 포화 여부, 적분 상태, 목표가 step인지 먼저 확인한다 \\
    Lab03 & 손끝 경로가 목표를 따라가고, 조건수가 급격히 나빠지지 않으면 관절 속도와 effort가 완만하다 & 손끝 오차는 작아 보이는데 관절 속도나 effort가 튄다 & 조작성, 조건수, 관절 속도, DLS 감쇠 설정을 함께 본다. 여기서 DLS 감쇠는 물리 감쇠가 아니라 역문제의 수치 완화이다 \\
    Lab04 & 접촉 전에는 침투 깊이와 \path|force_virtual_0|이 0에 가깝고, 벽을 넘은 뒤에만 침투, 계산된 벽 힘, 목표 후퇴가 함께 나타난다 & 접촉 전부터 벽 힘이 나오거나, 침투는 줄었는데 effort 지표가 커진다 & 벽 위치, 침투 깊이 부호, \path|force_virtual_0|과 \path|tau_cmd_*|/\path|current_proxy_*| effort proxy의 의미 차이를 먼저 본다 \\
    \bottomrule
  \end{tabularx}
\end{table}

표~\ref{tab:lab-first-plot-check}는 처음 플롯을 열었을 때의 확인 순서를 모은 것이다.
기대 모습은 정상 흐름의 기준선이고, 첫 이상 신호는 다음에 확인할 단서이다.
따라서 같은 감쇠라는 단어라도 Lab01의 물리 감쇠, Lab03의 DLS 수치 감쇠, Lab04의 벽 감쇠는 서로 다른 그래프에서 확인해야 한다.
Lab03의 DLS 감쇠는 역기구학 계산을 완만하게 만드는 수치 완화이고, Lab04의 벽 힘은 정의한 가상 벽 식에서 계산한 교육용 신호이다.

\subsection{Lab01: 질량-스프링-댐퍼}

Lab01은 식 \eqref{eq:msd}의 1차원 질량-스프링-댐퍼 시스템을 MuJoCo 슬라이드 관절로 구현한다.
학생은 질량, 강성, 감쇠, 초기 위치, 외력을 바꾸며 부족감쇠, 임계감쇠, 과감쇠 응답을 관찰한다.
이 랩은 임피던스 제어의 가장 작은 물리 블록이다.
기대 플롯은 위치, 속도, 가속도 추정값, 적용 힘, 운동에너지, 퍼텐셜 에너지이다.
예상되는 기본 그림은 위치가 평형점 쪽으로 움직이고, 부족감쇠에서는 몇 번 지나쳤다가 줄어들며, 과감쇠에서는 목표를 지나치지 않고 천천히 접근하는 모습이다.
강성 $k$를 키우면 같은 일정한 힘에서 최종 변위는 작아지고, 질량이 같다면 고유진동수가 커져 진동 주기가 짧아진다.
감쇠 $b$를 키우면 오버슈트와 에너지 교환의 진폭은 줄지만, 너무 크면 목표로 돌아오는 시간이 길어질 수 있다.
질량 $m$을 키우면 같은 힘에 대한 초기 가속도가 작아져 응답이 더 무겁고 느리게 보인다.
초심자는 최종 위치만 보고 감쇠나 질량 효과를 판단하기 쉽다.
하지만 감쇠와 질량의 영향은 주로 멈추기 전의 흔들림, 속도 변화, 에너지 변화에서 드러난다.

\subsection{Lab02: PID 제어}

Lab02는 같은 1차원 플랜트 위에 PID 제어를 올려, 물리 시스템과 제어기가 어떻게 만나는지 보여준다.
음의 피드백 부호가 맞고 목표가 정지해 있을 때, P 항은 위치 오차에 비례한 힘을 만들기 때문에 가상 스프링처럼 볼 수 있다.
D 항은 속도에 비례한 저항을 만들기 때문에 가상 댐퍼처럼 볼 수 있다.
I 항은 강성이나 감쇠가 아니라 오래 남아 있는 오차를 누적해 정상상태 오차를 줄이는 보정이며, 포화가 있으면 적분항이 과도하게 쌓이는 적분 포화(windup)를 만들 수 있다.
따라서 Lab02는 단순한 제어기 실습이 아니라, 임피던스 제어에서 강성과 감쇠를 이해하기 위한 중요한 다리이다.
이때 함께 볼 플롯은 목표 위치와 실제 위치, 위치 오차, 제어 입력, 그리고 I 항이 있는 경우 적분 상태이다.
플롯에서는 P, D, I 변화가 오차, 제어 입력, 오버슈트, 적분 상태에 어떻게 나타나는지 확인한다.
제어 입력 그래프가 위아래 한계에서 평평하게 잘리면 이득이 더 필요한 상황이 아니라 포화가 응답을 지배하는 상황일 수 있다.
이때는 위치 오차와 적분 상태를 함께 보아야 한다.

\subsection{Lab03: 2자유도 매니퓰레이터}

Lab01과 Lab02에서는 한 방향 위치만 봤지만, 로봇팔에서는 여러 관절이 함께 움직여 하나의 손끝 위치를 만든다.
Lab03은 2자유도 평면 매니퓰레이터를 사용해 절~\ref{sec:robotics-foundations}에서 본 관절공간과 작업공간의 차이를 플롯으로 보여준다.
임피던스 제어 전에 이 기구학을 먼저 보는 이유는 손끝의 강성, 감쇠, 힘도 결국 자코비안과 관절 움직임을 통해 구현되기 때문이다.
손끝에서는 단순한 스프링처럼 보이는 명령도 관절공간에서는 자세, 특이점, 조건수에 따라 전혀 다른 크기의 관절 변화와 effort로 나타날 수 있다.
특히 upper/lower branch 비교는 같은 손끝 목표를 두고도 팔꿈치를 어느 쪽으로 접는지에 따라 관절공간의 이야기가 달라진다는 것을 보여준다.
따라서 이 비교에서는 손끝 경로가 서로 비슷한지만 보지 말고, \(q_1,q_2\) 궤적, 조작성, 조건수, effort 지표가 어느 구간에서 갈라지는지 함께 본다.
이것이 역기구학의 다중해를 실험 결과로 읽는 가장 쉬운 방법이다.
여기서 branch 비교와 특이점/DLS 비교는 서로 연결되지만 같은 실험은 아니다.
branch 비교는 여러 가능한 자세 중 어느 가지를 택했는지에 초점을 두고, 특이점/DLS 비교는 선택한 자세 근처에서 자코비안 역문제가 얼마나 예민해지는지에 초점을 둔다.
여기서 몇 가지 로봇 용어를 먼저 잡아두자.
자코비안은 작은 관절 움직임이 손끝 움직임으로 어떻게 바뀌는지를 나타내는 행렬이다.
특이점은 어떤 방향의 손끝 움직임을 만들기 어려워지는 자세이다.
조작성(manipulability)은 손끝을 여러 방향으로 얼마나 쉽게 움직일 수 있는지를 나타내는 지표이다.
조건수(condition number)는 자코비안 역문제가 얼마나 민감한지를 나타내는 숫자이고, 값이 커질수록 작은 손끝 목표가 큰 관절 변화로 바뀌기 쉽다.
감쇠 최소제곱(DLS)은 불안정한 역기구학 계산에 감쇠를 넣어 너무 큰 관절 변화가 나오지 않게 완화하는 방법이다.
여기서 $q_1$, $q_2$는 첫 번째와 두 번째 관절각이고, $l_1$, $l_2$는 두 링크의 길이이다.
정기구학은
\begin{align}
  x &= l_1 \cos q_1 + l_2 \cos(q_1 + q_2), \\
  y &= l_1 \sin q_1 + l_2 \sin(q_1 + q_2)
\end{align}
로 표현된다.
자코비안은
\begin{equation}
  \bm{J}(\bm{q}) =
  \begin{bmatrix}
  -l_1\sin q_1 - l_2\sin(q_1+q_2) & -l_2\sin(q_1+q_2) \\
   l_1\cos q_1 + l_2\cos(q_1+q_2) &  l_2\cos(q_1+q_2)
  \end{bmatrix}.
\end{equation}
이 자코비안을 이용하면 작업공간 오차를 관절 토크로 바꾸는 자코비안 전치 PD 제어를 구현할 수 있다.
또한 2자유도 평면 팔에서는 $\det(\bm{J})=l_1l_2\sin q_2$이므로, 특이점 근처에서는 부호 있는 $\det(\bm{J})$가 단순히 작아진다기보다 $|\det(\bm{J})|$ 또는 조작성
\begin{equation}
  w=\sqrt{\det(\bm{J}\bm{J}^{\mathsf{T}})}=|\det(\bm{J})|
\end{equation}
이 0에 가까워지고, 조건수는 커지며, 역기구학 문제가 불안정해진다.
$w$가 0에 가까울수록 손끝을 특정 방향으로 움직이기 위해 관절이 과하게 움직여야 한다.
이때 감쇠 최소제곱 계열 방법은 특이점 근처의 수치적 폭주를 줄이는 데 유용하다 \cite{buss2005sdls,schinstock1994robustik}.
직관적으로는 팔을 완전히 편 자세를 생각하면 된다.
팔이 거의 일직선이면 손끝을 팔 방향으로 조금 더 뻗는 동작은 매우 어려워진다.
수학적으로는 자코비안이 어떤 방향의 작은 손끝 움직임을 충분히 만들지 못하는 상태이고, 이때 역기구학은 작은 목표 변화에도 큰 관절 변화로 반응할 수 있다.
DLS는 이 문제를 완전히 없애는 만능 해법이 아니라, 너무 큰 관절 속도나 수치적 폭주를 완화하기 위해 역문제에 감쇠를 넣는 방법이다.
Lab03 플롯에서는 손끝 경로만 보지 말고 관절 속도, 토크 명령 또는 전류 프록시 형태의 actuator effort 지표, 조작성, 조건수를 함께 본다.
특이점 근처에서는 손끝 위치 오차가 작아 보여도 관절 속도나 토크 요구가 갑자기 커질 수 있다.
특이점에 가까워지는 플롯에서는 조작성은 골짜기처럼 낮아지고 조건수는 봉우리처럼 커진다.
DLS 감쇠를 키운 뒤 관절 속도 피크가 낮아지는 대신 손끝 오차가 더 오래 남는다면, 이는 오류가 아니라 정확도와 안정성 사이의 의도된 절충으로 읽는다.
조건수 봉우리가 사라지지 않아도 이상한 일이 아니다.
DLS는 나쁜 자세를 갑자기 좋은 자세로 바꾸는 방법이 아니라, 그 자세에서 역기구학이 너무 큰 관절 속도 명령을 요구하지 않도록 완충하는 계산 방법이다.
따라서 Lab03 DLS 플롯은 조작성/조건수, 관절 속도 피크, actuator effort 지표, task error가 같은 시간대에서 어떻게 함께 변하는지 묶어서 읽는다.
또 같은 손끝 경로라도 목표가 시간에 따라 얼마나 빨리 움직였는지에 따라 관절 속도와 effort 피크가 달라질 수 있다.
따라서 path를 잘 따라갔는지와 trajectory가 너무 급했는지는 따로 읽어야 한다.
속도 플롯이 갑자기 꺾이거나 가속도 추정이 크게 튀면, 사다리꼴 속도 프로파일의 전환점이나 급한 목표 변경이 관절공간 요구를 키운 흔적으로 볼 수 있다.
Lab03의 오래된 1차원 궤적 프로파일 설정은 이 구분을 직접 확인하는 가장 작은 체크포인트이다.
\path|configs/lab03_2dof/step.yaml|, \path|configs/lab03_2dof/trapezoidal.yaml|, \path|configs/lab03_2dof/minimum_jerk.yaml|, \path|configs/lab03_2dof/s_curve.yaml|은 비슷한 이동 목표를 서로 다른 시간표로 읽게 해 준다.
이때 손끝 또는 관절 경로 오차만 보지 말고, 목표/실제 속도, 가속도, jerk 또는 jerk 추정치가 저장되어 있으면 그 값, 그리고 \texttt{tau\_cmd}나 \texttt{current\_proxy} 형태의 effort 피크가 같은 전환 시각에 함께 커지는지 확인한다.
단, 이상적인 사다리꼴 속도 프로파일의 전환점은 수학적으로 가속도가 불연속인 사건이다.
샘플된 로그에서 특정 열의 jerk 값이 0으로 보이더라도, 그것만으로 전환 충격이 물리적으로 없다고 해석하면 안 된다.

\begin{table}[!htbp]
  \centering
  \caption{경로와 궤적 구분을 랩에서 확인하는 안전한 읽기}
  \label{tab:trajectory-profile-lab-checkpoint}
  \small
  \setlength{\tabcolsep}{4pt}
  \renewcommand{\arraystretch}{1.12}
  \begin{tabularx}{\linewidth}{@{}>{\raggedright\arraybackslash}p{0.17\linewidth}>{\raggedright\arraybackslash}p{0.37\linewidth}>{\raggedright\arraybackslash}X@{}}
    \toprule
    랩 & 비교할 설정 & 안전한 해석 \\
    \midrule
    Lab03 프로파일 & \path|step.yaml|, \path|trapezoidal.yaml|, \path|minimum_jerk.yaml|, \path|s_curve.yaml| & 같은 목표라도 시간표가 달라지면 속도, 가속도, jerk, effort 피크가 달라진다. S-curve가 항상 더 좋다는 뜻은 아니며 제한 조건과 추종 오차를 같이 본다 \\
    Lab03 DLS/속도 제한 & \path|condition_aware_dls_slow_command_2dof.yaml|, \path|condition_aware_dls_fast_command_2dof.yaml|, joint-speed-limit 비교 설정 & 빠른 명령은 손끝 목표가 비슷해도 DLS 속도 제한, 관절 속도, task error의 절충을 더 크게 드러낼 수 있다 \\
    Lab04 벽 접근 & \path|wall_slow_approach.yaml|, \path|wall_fast_approach.yaml| & 같은 벽 강성/감쇠에서도 접근 시간과 속도가 달라지면 계산된 가상 벽 감쇠력과 effort 지표가 달라진다. 이는 S-curve 접촉 제어 성능 검증이 아니라 접근 속도 효과를 읽는 비교이다 \\
    \bottomrule
  \end{tabularx}
\end{table}

따라서 이 랩의 질문은 손끝이 목표에 갔는가만이 아니라, 그 목표를 만들기 위해 관절공간에서 얼마나 부담을 주었는가이다.

\subsection{Lab04: Panda 매니퓰레이터와 가상 벽}

Lab03에서 배운 손끝 목표를 관절 움직임으로 바꾸는 문제가 Lab04에서는 7자유도 로봇 모델 위로 올라간다.
Lab04는 Franka Emika Panda 모델로 7자유도 로봇팔의 손끝 병진 운동을 관찰하는 교육용 데모이다.
Franka 기반 MuJoCo 환경은 로봇 학습과 제어 실험에서 널리 사용되고 있으므로 \cite{xu2024frankamujocoenvs}, 이 절에서는 공개 모델 위에서 위치, 속도, 액추에이터 effort 지표, 계산된 가상 벽 신호가 앞의 임피던스 직관과 어떻게 맞물리는지 본다.
해석 범위는 손끝의 $x$, $y$, $z$ 위치에 둔다.
이때 7개 관절로 3차원 병진 위치를 맞추므로, 같은 손끝 위치를 만드는 관절 자세가 하나로 고정되지 않는다.
이 여유 자유도는 실제 로봇 제어에서는 null-space 자세 선택, 관절 한계 회피, 특이점 회피 같은 문제로 이어진다.
다만 현재 Lab04는 null-space 목적함수나 자세 최적화를 검증하는 실험이 아니라, 현재 자세에서 DLS 목표 오프셋이 어떤 관절 목표 변화와 effort 지표를 만드는지 읽는 교육용 장면이다.
현재 구현의 Menagerie 기반 Panda는 position 계열 actuator를 사용하므로 \texttt{ctrl}은 관절 목표 위치로 읽는다.
작업공간 도달(Cartesian reach)과 가상 벽(virtual wall) 데모는 작업공간 오차 또는 가상 벽 신호를 DLS 기반 관절 목표 오프셋으로 바꾸고, \texttt{tau\_cmd\_*}와 \texttt{current\_proxy\_*}는 MuJoCo \texttt{data.actuator\_force}에서 얻은 effort 기록으로 읽는다.
따라서 이 절은 6차원 자세 임피던스나 운영공간 역동역학 검증이 아니라, position actuator와 DLS 목표 오프셋 위에서 병진 응답을 읽는 훈련으로 한정한다.
즉, 손끝 위치 실험이지 실제 토크 기반 전체 임피던스 제어 검증은 아니다.
이 절에서 목표 후퇴(retreat)는 벽 힘이 커질 때 목표점을 벽 바깥쪽으로 물러나게 하여 침투와 힘을 줄이는 동작을 뜻한다.

표~\ref{tab:impedance-implementation-ladder}의 구현 사다리에서 말한 네 번째 층위가 바로 이 Lab04 데모이다.
신호 흐름은 다음 순서로 읽으면 된다.
\begin{enumerate}
  \item 끝단 위치와 속도에서 침투 깊이와 접근 속도를 계산한다.
  \item 그 값으로 \texttt{force\_virtual\_0}, \texttt{force\_virtual\_spring\_0}, \texttt{force\_virtual\_damping\_0}을 만든다.
  \item 이 힘 신호를 토크 명령으로 보내지 않고 목표점을 벽 바깥쪽으로 후퇴시킨다.
  \item DLS가 후퇴한 목표를 관절 목표 위치 오프셋으로 바꾼다.
  \item position actuator가 \texttt{ctrl}에 적힌 관절 목표 위치를 따라간다.
\end{enumerate}
즉, Lab04 플롯의 힘 이름들은 실제 힘을 측정했다는 뜻이 아니라, 계산된 벽 힘이 목표 후퇴와 actuator effort 지표로 이어지는 과정을 남긴 기록이다.

Lab04 신호는 다음처럼 나누어 읽으면 덜 헷갈린다.

\begin{table}[!htbp]
  \centering
  \caption{Lab04 주요 신호의 해석 범위}
  \small
  \setlength{\tabcolsep}{4pt}
  \renewcommand{\arraystretch}{1.12}
  \begin{tabularx}{\linewidth}{@{}>{\raggedright\arraybackslash}p{0.26\linewidth}>{\raggedright\arraybackslash}p{0.32\linewidth}>{\raggedright\arraybackslash}X@{}}
    \toprule
    신호 & 이 데모에서의 의미 & 해석할 때 구분할 점 \\
    \midrule
    \texttt{force\_virtual\_0} 및 spring/damping 성분 & 가상 벽 강성, 침투 깊이, 접근 속도로 계산한 벽 모델 출력 & 실측 접촉력이나 MuJoCo 접촉 제약력과 구분한다 \\
    MuJoCo \texttt{data.actuator\_force} & position actuator가 목표 관절 위치를 만들기 위해 낸 actuator force의 원자료 & CSV 열 이름으로 직접 저장되기보다 \texttt{tau\_cmd\_*}, \texttt{current\_proxy\_*}로 기록된다 \\
    \texttt{tau\_cmd\_*}, \texttt{current\_proxy\_*} & 다른 랩과 플롯 형식을 맞추기 위해 기록한 effort proxy & Lab04에서는 운영공간 토크 명령이나 실측 모터 전류가 아니다 \\
    \texttt{ctrl} & 관절 목표 위치 명령 & 토크 명령이 아니라 위치 계열 actuator 입력이다 \\
    \texttt{wall\_retreat\_cm}, \texttt{target\_wall\_gap\_m} & 계산된 가상 벽 신호에 따른 목표 후퇴량과 목표-벽 간격 & 실제 침투 깊이와 같지 않다. 제어기가 바꾼 평형점의 흔적으로 읽는다 \\
    \bottomrule
  \end{tabularx}
\end{table}

Lab04의 핵심 데모는 다음과 같다.
\begin{itemize}
  \item 기준 자세 유지: 로봇이 안정적으로 기준 자세를 유지하는지 확인한다.
  \item 관절공간 추종: 관절공간 목표 궤적과 실제 관절각의 차이를 관찰한다.
  \item 작업공간 도달: 끝단 위치 오차와 자코비안 기반 목표 보정을 확인한다.
  \item 가상 벽: 그림~\ref{fig:virtual-wall-comparison}의 관계를 바탕으로, 정의해 둔 벽 평면을 얼마나 넘었는지, 계산된 가상 벽 힘과 목표 후퇴 동작이 어떻게 달라지는지 비교한다.
\end{itemize}

가상 벽 플롯에서는 끝단 위치, 벽 위치, 침투 깊이, 벽 힘, 목표 후퇴량 \texttt{wall\_retreat\_cm}, 액추에이터 effort 지표를 한 묶음으로 읽는다.
절~\ref{sec:robotics-foundations}의 경로와 궤적 구분을 가져오면, 벽을 향해 가는 공간 경로뿐 아니라 접촉 직전의 시간표도 함께 봐야 한다.
같은 목표 침투 깊이라도 빠르게 접근하면 속도항과 과도응답이 커지고, 가속도나 jerk가 큰 목표 변경은 effort 피크와 튕김을 더 크게 만들 수 있다.
여기서 벽 힘은 실측 접촉력이 아니라 앞 절의 가상 벽 식에서 계산한 값이다.
끝단이 벽 바깥에 있을 때 침투 깊이와 벽 힘은 거의 0이고, 벽 안쪽으로 들어가면 $\delta$가 양수가 된다.
로봇이 $+x$ 방향으로 벽 안쪽에 들어가면, 벽은 반대로 $-x$ 방향으로 밀어낸다고 약속한다.
따라서 이 부호 규약에서는 \texttt{force\_virtual\_0}이 음수 방향으로 내려가며, 크기를 비교할 때는 \(\lvert\text{\texttt{force\_virtual\_0}}\rvert\)를 본다.
시간 순서로는 접촉 전, 진입 중, 정지 후를 나누어 본다.
접촉 전에는 침투 깊이와 벽 힘이 0이다.
진입 중에는 침투가 생기고 접근 속도가 양수이면 스프링항과 감쇠항이 함께 커진다.
정지한 뒤에는 속도항이 사라지므로 벽 힘은 주로 강성-침투 깊이 직선으로 남는다.
벽 강성 $k_{\mathrm{wall}}$을 키우면 같은 침투 깊이에서 벽 힘의 기울기가 커진다.
침투 깊이는 같은 목표 궤적, 목표 후퇴 게인, 포화 한계, DLS 조건에서는 작아지는 경향을 기대할 수 있지만, 위치 액추에이터 추종 성능과 목표 후퇴 제한이 섞이면 비단조적으로 보일 수 있다.
이때 액추에이터 힘과 튕김 위험은 더 커질 수 있다.
벽 감쇠 $b_{\mathrm{wall}}$을 키우면 가상 벽 안쪽으로 들어가는 동안의 속도항이 커져 접촉 순간의 튕김과 진동이 줄 수 있지만, 정지한 뒤의 힘-침투 직선 기울기는 바꾸지 않는다.
유효 질량이 큰 자세나 방향에서는 같은 강성과 감쇠를 써도 더 느리고 무겁게 반응하며, 충분히 감쇠되지 않으면 접촉 후 흔들림이 더 오래 남을 수 있다.

\begin{table}[!htbp]
  \centering
  \caption{랩 결과를 처음 읽을 때 놓치기 쉬운 점}
  \label{tab:lab-result-misreadings}
  \small
  \renewcommand{\arraystretch}{1.12}
  \begin{tabularx}{\linewidth}{@{}p{0.34\linewidth}X@{}}
    \toprule
    흔한 첫 해석 & 함께 볼 신호 \\
    \midrule
    최종 위치 오차만 보면 충분하다 & 과도응답, 속도, 제어 입력, 힘, 에너지를 함께 봐야 강성, 감쇠, 질량 효과를 구분할 수 있다. \\
    감쇠를 키우면 접촉력이 항상 줄어든다 & 감쇠는 주로 움직이는 동안의 튕김과 진동을 줄이고, 정지 후 힘-침투 기울기는 주로 강성이 정한다. \\
    뷰어 움직임이 부드러우면 항상 안전하다 & 화면상 움직임이 부드러워도 액추에이터 힘 지표나 토크 요구가 클 수 있다. \\
    가상 벽 힘은 MuJoCo 접촉 해석기 힘과 같다 & 이 데모의 벽 힘은 정의한 스프링-댐퍼 벽 모델에서 계산한 교육용 신호이다. 실제 접촉 제약력과 구분한다. \\
    목표 후퇴량과 침투 깊이는 같다 & 목표 후퇴는 의도적으로 움직인 평형점 위치이고, 침투 깊이는 실제 끝단이 정의된 벽 평면을 넘어간 양이다. \\
    \bottomrule
  \end{tabularx}
\end{table}

표~\ref{tab:lab-plot-symptom-guide}에서는 플롯이 예상과 다를 때 먼저 볼 신호를 정리한다.
여기서 중요한 것은 결과를 한 문장으로 단정하기보다, 어떤 신호가 어느 파라미터와 함께 움직였는지 확인하는 것이다.

\begin{table}[!htbp]
  \centering
  \caption{플롯 증상에서 시작하는 랩 결과 읽기}
  \label{tab:lab-plot-symptom-guide}
  \small
  \renewcommand{\arraystretch}{1.12}
  \begin{tabularx}{\linewidth}{@{}>{\raggedright\arraybackslash}p{0.27\linewidth}>{\raggedright\arraybackslash}p{0.27\linewidth}>{\raggedright\arraybackslash}X@{}}
    \toprule
    플롯에서 보이는 모습 & 먼저 볼 신호 & 가능한 해석 \\
    \midrule
    최종 변위가 예상보다 크다 & 외력, 강성, 정상상태 도달 여부 & 강성이 낮거나 아직 과도응답 중일 수 있다. 강성을 키웠는데도 줄지 않으면 부호와 입력 크기를 확인한다. \\
    짧은 주기로 심하게 흔들린다 & 감쇠비, 속도, 에너지, 제어 입력 & 강성이 커졌거나 감쇠가 부족할 수 있다. 포화가 있으면 단순 감쇠 증가만으로 해결되지 않는다. \\
    제어 입력이 한계에서 평평하게 잘린다 & actuator effort, 위치 오차, 적분 상태 & 이득 부족이 아니라 포화가 응답을 지배할 수 있다. I 항이 있으면 windup을 함께 의심한다. \\
    손끝 오차는 작은데 관절 속도나 effort가 튄다 & 조작성, 조건수, 관절 속도 & 특이점이나 관절 제한 근처에서 작은 작업공간 목표가 큰 관절공간 요구로 바뀐 것일 수 있다. \\
    침투 깊이는 줄었지만 벽 힘이나 effort가 커진다 & 벽 강성, 침투 깊이, \path|tau_cmd_*|/\path|current_proxy_*| effort proxy & 더 단단한 벽 또는 큰 목표 후퇴가 작은 변위를 큰 힘 지표로 바꾼 것일 수 있다. \\
    가상 벽 힘이 접촉 전부터 0이 아니다 & 끝단 위치, 벽 위치, 침투 깊이 계산 & 벽 위치 기준, 부호, 침투 깊이 계산을 먼저 확인한다. \\
    \bottomrule
  \end{tabularx}
\end{table}

표~\ref{tab:lab-plot-symptom-guide}에서 한 걸음 더 나아가면, 예상과 다른 플롯은 다음 실험을 고르는 기록이 된다.
먼저 관찰한 신호를 적고, 그 신호가 어떤 파라미터와 함께 움직였는지 해석한 뒤, 다음 실행에서 하나만 바꾼다.
한 번에 하나만 바꾼다는 말은 독립 원인을 분리한다는 뜻이다.
감쇠비를 유지하기 위해 감쇠계수를 함께 재계산하거나 같은 정적 힘 목표를 유지하기 위해 목표 오프셋을 맞추는 일은 같은 가설 안의 종속 조정으로 둔다.

\begin{table}[!htbp]
  \centering
  \caption{예상과 다른 플롯에서 다음 실험으로 넘어가기}
  \label{tab:lab-failure-case-next-run}
  \small
  \setlength{\tabcolsep}{4pt}
  \renewcommand{\arraystretch}{1.12}
  \begin{tabularx}{\linewidth}{@{}>{\raggedright\arraybackslash}p{0.29\linewidth}>{\raggedright\arraybackslash}p{0.31\linewidth}>{\raggedright\arraybackslash}X@{}}
    \toprule
    관찰한 모습 & 이렇게 읽는다 & 다음에 하나만 바꿔 볼 것 \\
    \midrule
    최종 위치 오차나 침투 깊이가 예상보다 크고 effort 지표가 포화되지 않았다 & 너무 부드러운 설정이라 같은 힘에서 많이 밀릴 수 있다. Lab04에서는 목표 후퇴 게인이나 제한이 함께 영향을 줄 수 있다 & Lab01은 \texttt{stiffness}를 조금 올린다. Lab04 가상 벽은 같은 나머지 조건을 유지한 채 \texttt{virtual\_wall.stiffness}, 목표 후퇴량, 접근 속도 중 하나만 조정한다. 감쇠비를 유지하는 1축 모델이라면 감쇠계수도 함께 다시 계산한다. \\
    목표 도달 뒤나 벽 접촉 뒤 짧은 주기로 튕긴다 & 강성에 비해 감쇠가 부족한 과도응답일 수 있다 & 1D나 작업공간 데모에서는 $\zeta$ 또는 해당 damping 설정을 조금 올리고, 가상 벽 데모에서는 \texttt{virtual\_wall.damping}을 조금 올린다. 정지 후 힘-침투 기울기는 강성이 주로 정한다. \\
    제어 입력이나 effort 지표가 한계에서 평평하게 잘린다 & 이득 부족이 아니라 포화, 목표 오프셋 제한, 관절 목표 변화 제한이 응답을 지배할 수 있다 & 강성 계열 설정, 목표 속도, 목표 오프셋 중 하나를 낮추고, 더 큰 값을 넣기 전에 제한값을 확인한다. \\
    침투 깊이는 줄었지만 계산된 벽 힘이나 effort 지표가 더 커졌다 & 더 안전해진 것이 아니라, 작은 변위를 큰 힘으로 막는 더 단단한 설정이 되었을 수 있다 & \texttt{virtual\_wall.stiffness}, 목표 후퇴량, 접근 속도 중 하나만 조정한다. Lab04의 벽 힘은 결정론적 가상 벽 신호임을 함께 기록한다. \\
    가상 벽 힘이 접촉 전부터 0이 아니거나, $+x$ 방향 침투 중 부호가 예상과 다르다 & 벽 위치 기준, 침투 깊이 부호, 힘의 크기와 부호 있는 값을 섞어 읽었을 수 있다 & $\delta=\max(0,x_{\mathrm{ee}}-x_{\mathrm{wall}})$ 계산과 플롯 축을 먼저 확인한다. 크기 비교는 $|f_{\mathrm{wall}}|$로 하고, 부호 있는 힘은 방향 약속과 함께 본다. \\
    손끝 오차는 작은데 관절 속도나 effort 지표가 튄다 & 작업공간 오차는 작아 보여도 특이점이나 관절 제한 근처에서는 작은 손끝 움직임이 큰 관절공간 요구로 바뀔 수 있다 & DLS 감쇠를 한 단계 올리고 손끝 오차가 얼마나 더 오래 남는지 함께 본다. 이 감쇠는 물리 감쇠비가 아니라 역문제의 수치 완화 파라미터이다. \\
    \bottomrule
  \end{tabularx}
\end{table}

이 흐름은 1차원 스프링-댐퍼 직관에서 시작해, PID의 P/D 항을 거쳐, 2자유도 작업공간 제어와 7자유도 매니퓰레이터의 임피던스적 상호작용까지 이어진다.
각 랩은 같은 순서로 읽는다.
설정값이 바뀌고, 시간응답이 달라지고, 그 차이가 힘이나 토크 지표로 다시 나타난다.
Lab01/02에서는 설정값 변화가 질량-스프링-댐퍼 관계에 비교적 직접 대응되고, Lab03의 torque-actuated 2DOF에서는 토크 명령을 함께 볼 수 있다.
반면 Lab04에서는 같은 직관을 position actuator 기반의 관절 목표 위치, DLS 목표 오프셋, 계산된 가상 벽 힘, 그리고 MuJoCo \texttt{data.actuator\_force}에서 나온 \texttt{tau\_cmd\_*}/\texttt{current\_proxy\_*} effort 지표로 읽어야 한다.

